\documentclass{jsarticle}
\usepackage{amssymb}
\usepackage{amsthm}
\usepackage{amsmath}
\usepackage{comment}
%定理環境 thmはあまり使わずpropをなるべく使う。
%主張は基本的に証明の中で使い、そのため番号を付けない。
%注意環境を付けてみた。
\newtheorem{thm}{定理}
\newtheorem{prop}[thm]{命題}
\newtheorem{cor}[thm]{系}
\newtheorem{lem}[thm]{補題}
\newtheorem{df}[thm]{定義}
\newtheorem{exam}[thm]{例}
\newtheorem{fact}[thm]{事実}
\newtheorem*{claim}{主張}
\newtheorem*{cau}{注意}
\renewcommand{\proofname}{証明}
%矢印たち。
%\toは\rightarrowと同じ。\getsは\leftarrowと同じ。同値は\iff
\newcommand{\To}{\Longrightarrow}
\newcommand{\Gets}{\LongLeftarrow}
%黒板太字
\newcommand{\nn}{\mathbb{N}}
\newcommand{\zz}{\mathbb{Z}}
\newcommand{\qq}{\mathbb{Q}}
\newcommand{\rr}{\mathbb{R}}
\newcommand{\cc}{\mathbb{C}}
%無限大
\newcommand{\mgn}{\infty}
%差集合は\setminusで統一する。等号の否定は\neq
%真部分集合は\subsetneqq　偏微分は\partial
%ディスプレイスタイル。\dfracでディスプレイ分数
\newcommand{\dpst}{\displaystyle}

\newcommand{\kk}{K^{\times}}
\newcommand{\ord}{\text{{\rm ord}}}
\DeclareMathOperator{\lcm}{lcm}
\title{体の乗法群の有限部分群について}
\author{電波通信}
\date{\today}
			\begin{document}
\maketitle

この文章では体$K$の乗法群$K^{\times}$の有限部分群が巡回群であることを示す。この文章ではわかりやすさのために自然数の整除関係を$\preceq$と書くことにする。つまり$n$が$m$の約数の時に$n\preceq m$と書くのである。$\le$と$\preceq$の違いに注意せよ。また、この文書において体とは可換体を指すことに注意せよ。

	\begin{comment}

	\begin{fact}[ラグランジュの定理]
	有限群$G$とその部分群$H$について
	\[
	|H|\preceq |G|
	\]
	となる。
	\end{fact}

	群$G$の元$a$の位数を$\ord(a)$と書くことにする。$a$の位数とは$a^n=e$となる最小の$n$であ			ることを一応書き記しておく。$e$は単位元である。
	また$a^m=e$なら$\ord(a)\preceq m$となることも思い出されたい。
	\begin{cor}
	有限群$G$の元の位数は$G$の約数。
	\end{cor}

	\end{comment}

\begin{df}群$G$の元$a$の位数$\ord(a)$を
$a^n=e$となる最小の$n$として定義する。このような$n$が存在しない場合には$\ord(a)=\mgn$とする。
\end{df}


\begin{prop}
群$G$の元$a$について$\ord(a)=n=xy$とすると
\[
\ord(a^x)=y
\]
となる。
\end{prop}
\begin{proof}
まず、$(a^{x})^y=a^{xy}=e$なので、$y\preceq \ord(a^x)$である。
そして、$a^{x\cdot\ord(a^x)}=(a^{x})^{\ord(a^x)}=e$なので$x\cdot\ord(a^x)\preceq n$となる。以上から$y\preceq \ord(a^x)\preceq y$なので$y=\ord(a^x)$である。
\end{proof}

準備のために次の整数に関する簡単な補題を証明する。
\begin{lem}
自然数$n,m$に対して$x,y,s,t$が存在して$n=sx$、$m=ty$、
$\gcd(x,y)=1$、$xy=\lcm(n,m)$を充す。
\end{lem}
\begin{proof}
自然数$n,m$に対して
$k$個の素数$p_1,p_2,\dots, p_k$と
$a_1, a_2, \dots, a_k\in \zz_{\ge0}$と
$b_1, b_2, \dots, b_k \in \zz_{\ge 0}$が存在して
\[
n=p_1^{a_1}p_{2}^{a_2}\cdots p_{k}^{a_k}
\]
\[
m=p_1^{b_1}p_{2}^{b_2}\cdots p_{k}^{b_k}
\]
と書ける。
このとき
\[
\lcm(n,m)=p_1^{\max(a_1,b_1)}p_{2}^{\max(a_2,b_2)}\cdots p_{k}^{\max(a_k,b_k)}
\]
である。
\begin{align*}
c_i=
	\begin{cases}
	a_i & \text{$a_i=\max\{a_i,b_i\}$のとき}\\
	0 & \text{$その他$}
	\end{cases}
\end{align*}
とし
\[
d_i=c_i-x_i
\]
と定義する。そして
\[
e_i=
	\begin{cases}
	b_i & \text{$b_i>a_i$のとき}\\
	0 & \text{$その他$}
	\end{cases}
\]
で
\[
f_i=b_i-z_i
\]
とする。
そして
\begin{align*}
x=p_1^{c_1}\cdots p_1^{c_2}\cdots p_k^{c_k}\\
s=p_1^{d_1}\cdots p_2^{d_2}\cdots p_k^{d_k}\\
y=p_1^{e_1}\cdots p_2^{e_2}\cdots p_k^{e_k}\\
t=p_1^{f_1}\cdots p_2^{f_2}\cdots p_k^{f_k}
\end{align*}
と定義すると、明らかに$n=sx$、$m=ty$、
$\gcd(x,y)=1$を満たし、また、定義から
\[
xy=(p_1^{c_1}\cdots p_1^{c_2}\cdots p_k^{c_k})(p_1^{e_1}\cdots p_2^{e_2}\cdots p_k^{e_k})=p_1^{\max(a_1,b_1)}p_{2}^{\max(a_2,b_2)}\cdots p_{k}^{\max(a_k,b_k)}=\lcm(n,m)
\]
を充す。
\end{proof}

\begin{thm}\label{aaa}
群$G$の元$a,b$は$ab=ba$を充し、$\ord(a)=n,\ \ord(b)=m$とし、$l=\lcm (n,m)$とする。この時$c\in G$が存在して 
\[
\ord(c)=l
\]
となる。
\end{thm}
\begin{proof}
まず$\gcd(n,m)=1$の時、$c=ab$とおき、$r=\ord(c)$とおくと$l=nm$よりと$a$と$b$の可換性より
\[
c^l=a^lb^l=ee=e
\]
なので$r\preceq l$でありまた
\[
e=c^{nr}=a^{nr}b^{nr}=eb^{nr}=b^{nr}
\]
より$m\preceq nr$ところで$n,m$は互いに素なので$m\preceq r$となる。同様に$n\preceq r$もわかるので結局$l\preceq r$（最小公倍数の性質である。）よって$\ord(c)=l$\\ 
一般の場合は$n=sx,\ m=ty,\ g.c.d.(x,y)=1\ xy=l$と表せるので$a^s$と$b^t$を上の場合に当てはめれば結論が得られる。
\end{proof}
\begin{cau}
また非可換な二元については定理\ref{aaa}は成り立たないことに注意せよ。例えば三次の対称群などを考えよ。
\end{cau}

定理\ref{aaa}をアーベル群に適用すると次のことがわかる。

\begin{cor}
有限可換群$G$の元の$\le$に関する最大値を$N$とすると、これは$G$の元が取りうる位数の中で$\preceq$に関しても最大値となっている。つまり有限可換群$G$のある元$x$が存在して
\[
\forall a\in G\ \ord(a)\preceq \ord(x)
\]
となるということである。
\end{cor}
\begin{proof}上の定理を$G$の元に次々に適用すればわかる。

\end{proof}

\begin{thm}
体$K$の乗法群$\kk$の有限部分群は巡回群
\end{thm}

\begin{proof}
$G$を$\kk$の有限部分群とする。$G$の最大位数を持つ元を$a$とし、$\ord(a)=N$と置く。$G$の任意の元$s$について$\ord(s)\preceq N$なので
\[
s^N=1
\]
となる。つまり$G$の任意の元は$K$上の多項式$x^N-1$の根である。$x^N-1$の根の個数は高々$N$個なので
\[
|G|\le N
\]
である。ところで$N\le |G|$でもあるから$N=|G|$よって$a$は$G$を生成する。つまり$G$は巡回群。
\end{proof}
































	
%\begin{thebibliography}{9}
%\bibitem{1}

%\end{thebibliography}




			\end{document}



\begin{comment}

[字体の変更]

{\rm hogehoge}と使う。

\rm ローマン
\it イタリック
\sl 斜体

[supp]

\newcommand{\supp}{\text{{\rm supp}}}

[中央揃えの改行]

	\begin{eqnarray*}
		hoge1&=&hogehoge
		hoge2&=&hogehofe

	\end{eqnarray*}

&&の部分で揃えられる。






[場合分け]

	\begin{cases}
		hoge1 & \text{状況１}\\
		hoge2 & \text{状況２}\\
		・
		・
		・
	\end{cases}



\end{comment}





